\section{Conclusion}

Systems that serve agentic LLM programs have been trying to recover, at the
serving layer, structure that the program already contains. When the LLM call is a
language construct rather than a string handed across an API boundary, that
structure is available to the compiler for free: the call topology, the provenance
of every argument, and the invariant portion of every prompt fall out of a static
pass with no annotation and no training. \sysname{} turns that knowledge into idle-time
work, constructing the longest knowable prefix of a future call during the windows
that tool calls and retrieval already open, and verifying it byte-for-byte against
byLLM's own rendering so that a wrong guess costs only a cache miss.

Two verification case studies on Qwen3-4B on a single RTX 3090 show the effect and
its boundary. On a supervisor-plus-four-subagents workflow, \(\Sigma\)TTFT falls
from 469.9\,ms to 265.1\,ms, capturing 72\% of the headroom to a measured
fully-cached Oracle, with the late, argument-heavy call sites improving by up to
5.07$\times$ and the call sites downstream of an unresolved branch improving not at
all. On a ReAct loop, the terminal call site improves 1.71$\times$ while the call
site with no preceding idle window is unchanged. In both cases the gain lands
exactly where provenance says a value is available before its consumer runs, which
is the strongest available evidence that the mechanism is doing what the design
claims rather than something incidental. Extending the evaluation across more agent
applications, model sizes, and concurrency levels is ongoing work.
